% 原书物理页=印刷页14--16；2.2 全节.
\section{半直积与扩张}\label{sec:fg-2-2}
\begin{numpara}\label{np:fg-2-2-1}
设 $A$ 与 $C$ 是两个群. 称 $C$ 在 $A$ 上的一个（左）作用与 $A$ 的群结构相容，如果相应的同态 $\tau:C\longrightarrow\mathrm{S}_A$ 在 $\operatorname{Aut}(A)$ 中取值. 换言之，若把这个作用记作 ${}^{c}a=\tau(c)(a)$，则要求对一切 $c,a_1,a_2$ 有 ${}^{c}(a_1a_2)={}^{c}a_1.{}^{c}a_2$.
\end{numpara}
\begin{proposition}\label{prop:fg-2-2-2}
用上述记号，赋予内部运算
\[
(a,c),\,(a',c')\longrightarrow(a.{}^{c}a',\,c.c')
\]
的集合 $A\times C$ 是一个群，其单位元为 $(1,1)$（并且有 $(a,c)^{-1}=({}^{c^{-1}}a^{-1},\,c^{-1})$）.
\end{proposition}
该命题中定义的群称为 $C$ 关于 $A$ 的、相对于 $\tau:C\longrightarrow\operatorname{Aut}(A)$ 的\newterm{fg-semidirect-product}{半直积}{semi-direct product}，记作 $A\times_\tau C$.
\begin{numpara}[内射映射]\label{np:fg-2-2-3}
此时，映射 $i:A\longrightarrow A\times_\tau C:\ i(a)=(a,1)$ 与 $s:C\longrightarrow A\times_\tau C:\ s(c)=(1,c)$ 都是单的群同态，从而可以把 $A$ 与 $C$ 分别等同于 $A\times_\tau C$ 的子群. 在这一等同之下，$A$ 在 $G$ 中正规\footnote{译注：此处及下文两处的 $G$ 泛指半直积群 $A\times_\tau C$.}，并且 $G$ 在 $A$ 上的共轭作用在 $C$ 上诱导出的恰好是 $\tau$：${}^{c}a=cac^{-1}$. 反之，$C$ 一般在 $G$ 中不正规：它正规当且仅当该作用是平凡的（$\tau=\mathit{id}_A$）；此时半直积不是别的，正是群 $A$ 与 $C$ 的直积 $A\times C$. 映射 $\phi:A\times C\longrightarrow A\times_\tau C:\ (a,c)\mapsto ac$（即 $i(a)s(c)$）总是一个集合间的双射，但仅当 $\tau$ 平凡时才是群的同构.
\end{numpara}
\begin{numpara}[投影]\label{np:fg-2-2-4}
投影 $p:A\times_\tau C\longrightarrow C,\ (a,c)\mapsto c$ 是一个满的群同态，其核为 $A$，后者通过上文定义的内射 $i$ 等同于 $A\times_\tau C$ 的一个子群. 反之，投影 $A\times_\tau C\longrightarrow A,\ (a,c)\mapsto a$ 一般\emph{不是群同态}. 它是群同态当且仅当 $\tau$ 是平凡作用，此时其核为 $C$，后者通过上述内射 $s$ 等同于 $A\times_\tau C$ 的一个子群.
\end{numpara}
\begin{numpara}[例：二面体群]\label{np:fg-2-2-5}
设 $n$ 为整数，$n\geqslant 2$. \newterm{fg-dihedral-group}{二面体群}{dihedral group} $D_n$ 是 $\mathrm{O}_2(\mathbb{R})$ 中以原点为中心的正 $n$ 边形（当 $n=2$ 时是一条线段）的稳定子. 它与 $\mathrm{SO}_2(\mathbb{R})$ 的交是由转角为 $2\pi/n$ 的旋转 $r$ 生成的 $n$ 阶循环群 $C_n$. 该子群的指数整除 $[\mathrm{O}_2(\mathbb{R}):\mathrm{SO}_2(\mathbb{R})]=2$，而且等于 $2$，因为 $D_n$ 包含关于连接原点与某个顶点的直线的对称中的任意一个（设为 $s$）. 于是有
\[
r^n=1,\ s^2=1,\ srs=srs^{-1}=r^{-1},
\]
从而对任意整数 $k$ 有 $sr^{-k}s=r^k$. 这样，$D_n$ 是一个 $2n$ 阶的群，同构于半直积 $C_n\times_\tau\{\pm1\}$，其中 $\{\pm1\}$ 在 $C_n$ 上的作用由 $\tau(-1)(a):=a^{-1}$ 给出. $D_n$ 的 $2$ 阶元素是那 $n$ 个对称 $sr^k$，以及当 $n$ 为偶数时的旋转 $r^{n/2}$. 当 $n=2$ 时，这个积是直积，此时群 $D_2$ 同构于 $C_2\times C_2$，常记作 $\mathbf{V}$（Vierergruppe，四元群；它作为指数为 $3$ 的正规子群出现在 $\mathrm{A}_4$ 中，因此 $\mathrm{A}_4$ 可解）. 群 $D_3$ 同构于对称群 $\mathrm{S}_3$.
\end{numpara}
\begin{numpara}[扩张]\label{np:fg-2-2-6}
一个群态射序列
\[
A\ \stackrel{i}{\hookrightarrow}\ E\ \stackrel{p}{\twoheadrightarrow}\ C
\]
若其中 $i$ 是单射、$p$ 是满射，且 $\operatorname{Im}(i)=\operatorname{Ker}(p)$，则称之为一个\newterm{fg-short-exact-sequence}{短正合列}{short exact sequence}. 也称 $(E,i,p)$ 是 $C$ 关于 $A$ 的一个\newterm{fg-extension}{扩张}{extension}. 因此，$C$ 关于 $A$ 的半直积特别地是 $C$ 关于 $A$ 的一个扩张. 我们想知道：一个扩张在什么条件下是某个适当作用 $\tau$ 下的半直积，以及这个作用 $\tau$ 在何种程度上被完全确定. 为此，把任何一个使得 $\varphi\circ i=i'$ 且 $p'\circ\varphi=p$ 的群态射 $\varphi:E\longrightarrow E'$ 称为一个\newterm{fg-extension-morphism}{扩张态射}{morphism of extensions} $(E,i,p)\xrightarrow{\ \varphi\ }(E',i',p')$.
\end{numpara}
\begin{proposition}\label{prop:fg-2-2-7}
设 $(E,i,p)$ 是 $C$ 关于 $A$ 的一个扩张. 那么，$(E,i,p)$ 同构于某个半直积 $(A\times_\tau C,i,p)$ 当且仅当 $p$ 存在一个\newterm{fg-section}{截面}{section}，即一个满足 $p\circ s=\mathrm{Id}_E$\footnote{译注：原文如此印作 $p\circ s=\mathrm{Id}_E$；由 $s:C\longrightarrow E$ 与 $p:E\longrightarrow C$ 知应为 $p\circ s=\mathrm{Id}_C$.} 的群态射 $s:C\longrightarrow E$. 此时称扩张 $(E,i,p)$ 是\newterm{fg-split-extension}{可裂的}{split}.
\end{proposition}
\begin{numpara}[例：四元数群]\label{np:fg-2-2-8}
考察 Hamilton 四元数体 $\mathbb{H}=\mathbb{R}\oplus\mathbb{R}i\oplus\mathbb{R}j\oplus\mathbb{R}k$，其乘法运算由 $\mathbb{R}$-双线性性并依照关系 $i^2=j^2=k^2=-1$，$ij=-ji=k$，$jk=-kj=i$，$ki=-ik=j$ 来定义（可以验证该运算确实是结合的）. \newterm{fg-quaternion-group}{四元数群}{quaternion group} $H_8$ 是乘法群 $\mathbb{H}^\times$ 的子群
\[
H_8=\{\pm1,\ \pm i,\ \pm j,\ \pm k\}.
\]
它的阶为 $8$，恰有三个 $4$ 阶循环子群 $\langle i\rangle$、$\langle j\rangle$、$\langle k\rangle$（与之相反，二面体群 $D_4$ 只有一个），并且只有一个 $2$ 阶元素（而 $D_4$ 则有 $5$ 个）. 由元素 $i$ 定义的 $C_4$ 到 $H_8$ 的内射赋予 $H_8$ 一个 $C_2$ 关于 $C_4$ 的扩张结构. 这个扩张不可裂，因为 $H_8$ 中唯一的 $2$ 阶元素是 $-1$. $\pm1$ 到 $H_8$ 的内射同样赋予 $H_8$ 一个（中心的、不可裂的）$\mathbf{V}$ 关于 $C_2$ 的扩张结构.
\end{numpara}
\begin{numpara}[1-上循环]\label{np:fg-2-2-9}
设 $C$ 是一个通过 $\tau$ 作用在另一个群 $A$ 上的群. 作用 $\tau$ 的一个\newterm{fg-1-cocycle}{1-上循环}{1-cocycle}是一个映射 $C\xrightarrow{\ \alpha\ }A$，使得
\[
\forall c_1,c_2\in C,\ \alpha(c_1c_2)=\alpha(c_1){}^{c_1}\alpha(c_2).
\]
\end{numpara}
\begin{proposition}\label{prop:fg-2-2-10}
设 $\tau$ 与 $\tau'$ 是 $C$ 在 $A$ 上的两个作用. 那么，相应的两个半直积作为 $C$ 关于 $A$ 的扩张是同构的，当且仅当存在作用 $\tau$ 的一个 1-上循环 $\alpha$，使得
\[
\forall c\in C,\ \tau'(c)=\tau(c)\circ\operatorname{Int}(\alpha(c)).
\]
\end{proposition}
\begin{numpara}[识别半直积]\label{np:fg-2-2-11}
反过来，设 $G$ 是一个群，$H,K$ 是 $G$ 的两个子群，使得映射 $\phi:H\times K\longrightarrow G:\ (h,k)\mapsto hk$ 是集合间的双射. 这特别地蕴含 $H\cap K=\{1\}$；而当 $G$ 有限时，它等价于这一条件与关系 $|G|=|H|.|K|$ 的合取. 若 $H$ 与 $K$ 交换，也就是说对一切 $(h,k)\in H\times K$ 都有 $kh=hk$，则映射 $\phi$ 是一个群同构（关于 $H\times K$ 上的乘积运算），从而 $G$ 是 $H$ 与 $K$ 的直积.

子群 $H$ 与 $K$ 交换当且仅当 $K$ 包含在 $H$ 的中心化子 $Z_G(H)$ 之中. 现在只假设 $K$ 包含在 $H$ 的正规化子 $N_G(H)$ 之中（由于 $HK=G$，这等价于要求 $H$ 在 $G$ 中正规）. 这时，$K$ 通过共轭作用在 $H$ 上：$(k,h)\mapsto\tau(k)(h):=khk^{-1}$，并且 $G$ 同构于 $K$ 关于 $H$ 的、相对于 $\tau=(\operatorname{Int}_{|K})_{|H}$ 的半直积 $H\times_\tau K$.
\end{numpara}
\begin{numpara}[例]\label{np:fg-2-2-12}
设 $k$ 是一个交换域. $\mathrm{GL}_n(k)$ 的上三角子群 $B_n$ 是由对角矩阵组成的子群 $T_n\simeq(k^\times)^n$ 关于严格上三角子群 $U_n$ 的半直积（$T_n$ 在 $U_n$ 上的作用由 $\mathrm{GL}_n(k)$ 中的共轭给出）.

$k$ 上仿射直线的仿射变换群 $\mathit{Aff}^1(k)$ 是由以 $0$ 为中心的位似组成的乘性子群（$\simeq k^\times$）关于由平移组成的加性子群（$\simeq k$）的半直积（$k^\times$ 通过 $\mathit{Aff}^1(k)$ 中的共轭作用在 $k$ 上，这个作用就是乘法本身）.
\end{numpara}
